%% kapa: exam=1 round=37 year=2026 period=1 form=A subject=civil_law subject_ko=민법 seq=18 number=18
%% source: Q-Net 감정평가사 기출문제 (cst00309) artlSeq=5258924 file=raw/1008/5258924/2026년 제37회 감정평가사 1차 1교시.pdf
%% conversion: kordoc 4.24.51 -> tools/gen1.py (auto); review=proofread-sonnet; flags=none
\begin{question}{18}
법률행위의 부관에 관한 설명으로 옳지 않은 것은? (다툼이 있으면 판례에 따름)
\begin{choices}
  \item 정지조건부 법률행위는 조건이 성취되었을 때 비로소 법률행위가 성립한다.
  \item 조건이 선량한 풍속 기타 사회질서에 위반한 것이면 그 법률행위는 무효이다.
  \item 법률행위에 조건이 붙어 있는지의 여부는 조건의 존재를 주장하는 자에게 증명책임이 있다.
  \item 조건이 법률행위 당시에 이미 성취할 수 없는 것인 경우, 그 조건이 해제조건이면 그 법률행위는 조건 없는 법률행위가 된다.
  \item 당사자가 불확정한 사실이 발생한 때를 이행기로 정한 경우에는 그 사실의 발생이 불가능하게 된 때에도 이행기가 도래한 것으로 보아야 한다.
\end{choices}
\end{question}
\answer{1}
